\chapter{多模态时序对齐与统一状态表示}

\section{多源异步采样问题分析}

无人机电磁超声检测的观测链条具有典型的异步多源特征：EMAT回波频率$f_Z \approx 40$ Hz，飞控和里程计反馈频率$f_P \approx 100$ Hz，RGBD视觉频率$f_V \approx 30$ Hz。三类传感器不仅采样频率悬殊，且在硬件层面存在固有延迟：视觉图像经历曝光（~15 ms）和深度计算（~10 ms），EMAT波形经历采样（~50 $\mu$s）和USB批量传输（~5 ms），飞控位姿经历EKF滤波延迟（~10 ms）和MAVROS串行转发（~5 ms）。若直接按原始时间戳拼接不同频率的数据，模型会将采样误差误认为物理变化，导致接触边界附近出现系统性误判。

\section{统一时间基准对齐方法}

为解决上述异步问题，理论框架首先建立统一时间基准，将不同模态映射到同一状态序列。设三类模态分别为不同时间戳上的观测序列$\{\mathbf{z}_t\}$、$\{\mathbf{p}_t\}$和$\{\mathbf{v}_t\}$，统一时间网格定义为等间隔的目标时间序列$\{t_k\}_{k=1}^K$，步长$\Delta t$取为三模态采样周期的最大公约数近似值（~10 ms）。

对任一模态$m$，在目标时刻$t_k$的观测值通过对局部时间窗口$[t_k - \frac{w}{2}, t_k + \frac{w}{2}]$内的原始样本进行加权插值获得：

\begin{equation}
\tilde{\mathbf{x}}_m(t_k) = \frac{\sum_{i \in \mathcal{N}(t_k, w)} \kappa(t_k, \tau_i) \cdot \mathbf{x}_m(\tau_i)}{\sum_{i} \kappa(t_k, \tau_i)}
\label{eq:weighted_interp}
\end{equation}

其中$\mathcal{N}(t_k, w)$为窗口$w$内的邻域索引集，$\kappa(t_k, \tau_i) = \exp(-\frac{|t_k - \tau_i|^2}{2\sigma^2})$为高斯核权函数，越接近目标时刻的原始样本赋予越高的插值权重。该步骤保留了原始传感器的采样特性（不做粗暴的零阶保持或最近邻重采样），同时避免了插值伪影造成的人为高频分量。

\section{模态独立编码器设计}

对齐后的三模态观测序列通过独立的编码器网络映射到统一维度的隐空间，以提取各模态中与接触状态判别最具相关性的特征：

\begin{equation}
\mathbf{h}_Z = f_Z(\tilde{\mathbf{Z}}; \theta_Z), \quad \mathbf{h}_P = f_P(\tilde{\mathbf{P}}; \theta_P), \quad \mathbf{h}_V = f_V(\tilde{\mathbf{V}}; \theta_V)
\label{eq:encoders}
\end{equation}

各编码器的结构设计遵循模态特性的差异：

\begin{itemize}
\item \textbf{超声编码器$f_Z$}：采用一维卷积神经网络结构，侧重捕获短时频谱分布、包络形状和主要峰值到达时间的变化。输入为原始波形切片（约1000采样点），经三层Conv1D（kernel size = 15, 31, 63）和最大池化层后，展平接入全连接层。
\item \textbf{位姿编码器$f_P$}：采用多层感知机结构，输入为位置误差向量$\mathbf{e}_p = \mathbf{p}_{\text{target}} - \mathbf{p}_{\text{current}}$、速度向量$\dot{\mathbf{p}}$和法向角度误差$\theta_n$的拼接。侧重捕获位姿的收敛/发散趋势和法向对齐度。
\item \textbf{视觉编码器$f_V$}：采用轻量级CNN（如MobileNetV3-Small前几层）提取目标ROI区域的空间特征，包括表面纹理梯度、目标边缘清晰度和深度图平滑度。
\end{itemize}

\section{位置编码与时序注入}

为保留时间顺序信息，在隐空间表示中引入正弦位置编码：

\begin{equation}
\mathrm{PE}_{(pos, 2i)} = \sin\left(\frac{pos}{10000^{2i/d}}\right), \quad \mathrm{PE}_{(pos, 2i+1)} = \cos\left(\frac{pos}{10000^{2i/d}}\right)
\label{eq:sinusoidal_pe}
\end{equation}

位置编码与各模态隐向量相加：$\mathbf{h}'_m = \mathbf{h}_m + \mathbf{PE}$，随后三类隐向量沿特征维度拼接形成统一状态表示$\mathbf{H} = [\mathbf{h}'_Z ;\ \mathbf{h}'_P ;\ \mathbf{h}'_V ] \in \mathbb{R}^{T \times d_{\text{total}}}$，作为后续Transformer层的输入。

\section{传感器延迟补偿}

工程实现中，多模态对齐还需处理传感器延迟。视觉图像经历曝光和深度计算，EMAT波形经历采样和USB传输，飞控位姿经历滤波和MAVROS转发。若这些延迟不被显式建模，数据集中的标签将发生系统性偏移。

本研究采用三种方式联合估计模态间延迟：时间戳对齐验证（记录各消息在ROS系统中的接收时间戳，与设备硬件时间戳交叉比对）、互相关峰值校正（对EMAT回波能量曲线与位姿法向残差曲线做滑动窗互相关，寻找最优滞后量）、固定延迟补偿（对已知延迟环节如USB传输和MAVROS转发施加固定前馈补偿）。三种方式联合确定的模态间最大偏移量控制在$\pm 15$ ms以内，满足接触状态估计的时间分辨率需求。

\section{本章小结}

本章针对多源异构传感器异步采样的核心问题，建立了统一的时序对齐方法论。通过局部时间窗口加权插值避免了粗暴重采样造成的信息损失，通过模态特定编码器将异构数据映射到统一隐空间，通过正弦位置编码保留时序顺序信息，并通过三种联合延迟补偿策略将模态间偏差控制在工程可接受范围。统一状态表示为后续章节的物理约束注意力机制提供了结构化的输入。

% ══════════════════════════════════════════════════════════════════
% 第五章 物理约束注意力机制
% ══════════════════════════════════════════════════════════════════